\section{Valutazione sperimentale}
\label{sec:valutazione}

\subsection{Ambiente e metodologia sperimentale}

La valutazione risponde a due domande. La prima riguarda la capacità del
livello Cloud di assorbire lo stesso workload riducendo ritardo e backlog
all'aumentare delle repliche. La seconda riguarda l'effetto di latenza e banda
disponibile sui due confini di rete Simulator--Edge ed Edge--Cloud. Le prove
sono eseguite sul deployment AWS di Fig.~\ref{fig:architecture}; ogni
configurazione viene ripetuta tre volte e viene considerata valida soltanto se
i controlli di qualità confermano la completezza della pipeline.

Il workload definitivo comprende 2\,106\,784 osservazioni distribuite tra 13
replay shard. Il replay usa un fattore di accelerazione pari a 20\,000; tale
valore è stato scelto mediante una calibrazione preliminare per esercitare il
livello Cloud senza introdurre intenzionalmente perdite nei componenti a
monte. La configurazione comune alle campagne è riassunta in
Tab.~\ref{tab:experiment-configuration}.

\begin{table}[t]
    \caption{Configurazione comune delle campagne sperimentali.}
    \label{tab:experiment-configuration}
    \centering
    \begin{tabular}{@{}lr@{}}
        \toprule
        Parametro & Valore \\
        \midrule
        Edge Site / Simulator & 13 \\
        Partizioni \code{edge-aggregates} & 6 \\
        Finestra Edge & 5 min \\
        Finestra Cloud & 15 min \\
        Fattore di accelerazione & 20\,000 \\
        Batch producer Edge & 10 \\
        Batch commit Cloud & 100 \\
        \bottomrule
    \end{tabular}
\end{table}

Per ogni run vengono conservati configurazione effettiva, log, metriche dei
container, serie temporali del consumer lag e aggregati finali PostgreSQL. La
riconciliazione verifica che gli eventi offerti dai Simulator coincidano con
quelli accettati dagli Edge, che non vi siano errori di pubblicazione o record
late, che tutte le finestre globali siano complete e che il lag finale sia
nullo. Le esecuzioni che non rispettano tali condizioni vengono documentate
ma escluse dal confronto prestazionale.

La metrica principale è la \emph{global emission latency}, cioè il ritardo con
cui una finestra globale viene emessa rispetto all'istante nominale previsto
dalla timeline accelerata. Indicando con $t_0$ l'avvio sincronizzato del
replay, con $e_0$ l'origine della timeline del dataset, con $f$ il fattore di
accelerazione e con $w_e$ la fine della finestra, la deadline nominale è

\begin{equation}
    d(w_e) = t_0 + \frac{w_e-e_0}{f}.
\end{equation}

La latenza è la differenza tra l'istante di emissione del Global Aggregator e
$d(w_e)$. Vengono riportati media, mediana, 95-esimo e 99-esimo percentile,
insieme al massimo consumer lag, all'utilizzo di CPU e memoria e al tempo
complessivo della run. Poiché il replay impone il ritmo di ingresso, il
makespan e il goodput sono usati soprattutto come controlli di capacità e
completezza; latenza e backlog evidenziano più direttamente il beneficio del
parallelismo Cloud.

\subsection{Scalabilità del livello Cloud}

La campagna usa un fixed workload e varia il numero di Cloud Worker tra
$W\in\{1,2,3,4,6\}$. Il topic mantiene sei partizioni in tutte le prove:
in questo modo cambia soltanto il numero di partizioni elaborate in parallelo,
mentre restano invariati dati, assegnamento degli Edge e risultati attesi. Sei
worker costituiscono il limite di parallelismo utile della configurazione;
repliche ulteriori rimarrebbero senza partizioni assegnate.

Il confronto utilizza le distribuzioni della global emission latency, il picco
e l'evoluzione del lag per partizione e l'utilizzo delle risorse. I valori
riportati nei grafici saranno la media delle tre repliche valide; la
variabilità tra run verrà mostrata esplicitamente, senza interpretare come
speedup la sola durata del replay, che è principalmente determinata dal rate
di ingresso.

% Risultati da inserire al termine della campagna:
% - tabella con run valide/scartate e motivazione;
% - grafico p50/p95/p99 della global emission latency per W;
% - grafico del massimo lag Cloud per W e, se utile, serie temporali W1/W6;
% - CPU dei worker e del Cloud Core;
% - discussione dello sbilanciamento fra le sei partizioni e del punto di
%   saturazione oltre il quale aggiungere worker non porta benefici.

\subsection{Impatto delle condizioni di rete}

Le condizioni di rete vengono emulate tramite Linux \code{tc-netem} nel
network namespace dei container interessati. Le regole sono selettive: sul
confine Simulator--Edge agiscono soltanto sul traffico MQTT verso il broker
associato; sul confine Edge--Cloud agiscono soltanto sulle connessioni Kafka.
Il traffico di orchestrazione e raccolta delle metriche non viene quindi
alterato.

La campagna usa tre Cloud Worker e confronta una baseline senza emulazione con
ritardi aggiuntivi di 50 e 100 ms e con limiti di banda applicati separatamente
ai due confini. Ogni profilo modifica una sola variabile alla volta. Oltre alla
latenza di emissione e al lag vengono controllati code, errori e perdite: una
riduzione apparente del carico ottenuta scartando telemetria non viene
interpretata come miglioramento prestazionale.

% Risultati da inserire al termine della campagna:
% - tabella con profilo, run valide, latenza e banda effettivamente applicate;
% - grafici latency/lag rispetto al ritardo e al rate sui due confini;
% - confronto delle code Simulator/Edge e dei tempi di pubblicazione Kafka;
% - distinzione fra degradazione controllata e perdita di completezza.

\subsection{Discussione dei risultati}

La discussione finale confronterà i benefici dell'aumento dei worker con i
limiti osservati: parallelismo massimo pari al numero di partizioni,
sbilanciamento residuo del carico, costo del coordinamento tramite Kafka e
possibile contesa del Cloud Core, che ospita broker e Global Aggregator. Le
conclusioni distingueranno inoltre la capacità di mantenere risultati completi
dalla sola velocità di elaborazione.

% Completare dopo le campagne con conclusioni quantitative, evitando di
% generalizzare oltre il deployment e il dataset effettivamente misurati.
